本附录集中说明报告口径、术语和不可据此推导的事项。其目的不是降低结论责任，而是确保事实、推导、情景和未披露项在投资行动中保持不同权重。

\section{关键口径定义}

\begin{exhibitbox}[术语与采用口径]
\small
\begin{tabularx}{\exhibitboxwidth}{L{3.7cm}L{6.0cm}X}
\toprule
\textbf{术语} & \textbf{本报告定义} & \textbf{不可替代项} \\
\midrule
2025年同一控制下重述口径 & 2025年报持续经营实际及其2024重述比较期 & 合并前legacy或交易备考 \\
当前股本 & 合并后75.25621288亿股 & 2025A EPS加权分母 \\
2025A EPS分母 & 公司披露加权平均63.29638855亿股 & 合并后期末股本 \\
在手订单 & 2025年末上市公司民船/海工聚合合同存量 & 中船集团整体、完整军品或无风险收入 \\
收入强度代理 & 核心收入除以交付DWT的分析代理 & 船型ASP或船东合同单价 \\
确认系数 & 相对2025收入确认强度的情景参数 & 公司披露履约进度 \\
Growth/Base & 用于驱动分析的业务拆分 & 公司法定新分部或独立估值单元 \\
中性偏多 & 方向偏正面但安全边际不足，持有等待验证 & 买入或收益承诺 \\
基础公允区间 & Base EPS对应11.0--12.5倍PE & 完整熊/牛情景范围 \\
市场锚 & 2026-07-22可观察现价 & 独立内在价值或技术趋势 \\
\bottomrule
\end{tabularx}
\end{exhibitbox}

\section{上市边界与零计价项目}

只有600150.SH是估值母体。七家主要船厂是consolidated components，其收入与利润已进入合并报表；中国动力是参股观察节点；沪东中华是集团姐妹单位、上市体外节点。任何子公司、参股或集团外资产都不能再次加入合并价值。

\begin{exhibitbox}[零基础估值信用清单]
\small
\begin{tabularx}{\exhibitboxwidth}{L{4.2cm}L{5.0cm}X}
\toprule
\textbf{项目} & \textbf{未计价原因} & \textbf{解除条件} \\
\midrule
未披露军品订单、收入和毛利 & 经济性与现金流无法审计 & 正式分部/合同披露及可复算盈利桥 \\
集团未分配订单 & 集团成绩不等于上市公司归属 & 600150公告明确合同与并表船厂 \\
沪东中华潜在注入 & 交易范围、价格、股本和盈利未公告 & 正式交易文件并整体重跑估值 \\
中远约500亿元项目增量 & 已在聚合订单内且含意向订单 & 不存在“解除去重”；只更新确定比例与兑现 \\
弱来源58.96/47.00元目标线索 & 无原始报告或可用正文 & 券商官方原件与完整假设 \\
大型邮轮独立溢价 & 单船利润、系列化经济性和系统价值不足 & 可审计订单、成本和现金证据 \\
集团LNG能力与排名 & 上市边界及公司份额未证明 & 上市公司订单/资产归属证据 \\
\bottomrule
\end{tabularx}
\end{exhibitbox}

零计价不代表这些项目永久没有价值，而是截至数据截止日没有足够证据进入基础、牛市或市场锚。未来若出现正式披露，应重新评估新增资产、对价、股本、少数股东和现金，而不是直接给一个主题溢价。

\section{单位与量纲}

报告金额除另有说明外为人民币亿元，股本为亿股，每股价值为人民币元。DWT衡量载重能力，CGT更接近建造工作量，二者不能互换；美元/CGT的全球混合价值量不能解释单船ASP。艘数、DWT、合同额、收入与利润是不同量纲，必须经交付和会计桥转换。

\begin{sourcequalitybox}[量纲控制]
在手合同额是存量，年度新签是流量，营业收入是履约确认，毛利是收入扣成本，归母净利还要考虑费用、税和少数股东。任何将这些层级直接相加、相除后解释为公司指引的做法都不属于本报告结论。
\end{sourcequalitybox}

\section{预测与目标价的不确定性}

2026--2028预测依赖交付、价格/结构、确认、毛利率、费用、减值、投资收益、税、少数股东和现金转化。船舶项目长周期且合同、成本与回款具有时点性，单季或半年结果可能偏离年度平均。H1预告为未经审计区间，不作机械年化。

目标价使用2027E周期正常化PE与有限市场锚，并不保证市场价格到达。基础价值、区间和情景会随正式数据变化；Bear/Base/Bull不是概率承诺，也不覆盖所有极端风险。+0.15\%是目标价相对现价的机械差值，不是预期收益保证。

\section{风险声明与使用限制}

本报告基于公开信息及其可复算推导，可能存在披露延迟、口径变动、模型误差与不可观察风险。地缘政治、制裁、贸易政策、客户融资、汇率、设备供应、质量事故、延期、撤单、重组整合和会计判断均可能使实际结果显著偏离预测。

本报告只用于投资研究和监测，不构成证券买卖建议、要约、承诺、受托管理或适当性判断。文中的“升级、持有、降低观察敞口”是研究动作语言，不是自动交易指令。投资者应独立判断并承担风险。

\begin{keyinsight}[复核原则]
任何人复用本报告时，必须同时保留：33.02元现价时点、75.2562亿股前瞻分母、2025同控重述口径、H1预告性质、4,674.51亿元订单边界、中远项目去重、唯一valuation parent和四类零计价。删除任一条件都会改变结论含义。
\end{keyinsight}

数据截止为2026年7月22日（北京时间）。在正式2026H1报告或新的公司公告出现后，本报告的盈利、估值、风险和动作均应重新审核，不应继续沿用旧目标价而只追加文字说明。
